\section{Related Work}

\subsection{Peptide--MHC Binding Prediction}

Peptide--MHC binding prediction is a foundational task in computational immunology. Established methods such as NetMHCpan \cite{reynisson2020netmhcpan}, MHCflurry \cite{odonnell2020mhcflurry}, and ACME \cite{hu2019acme} estimate binding affinity, binding rank, or peptide--MHC compatibility from peptide sequences and allele information. Pan-allelic predictors improve coverage by representing MHC molecules through pseudo-sequences or related allele features, thereby allowing information sharing across well-characterized and data-scarce alleles \cite{nielsen2007netmhcpan,reynisson2020netmhcpan}. These models are essential for filtering candidate epitopes, but binding is only one stage of antigen presentation. A peptide with favorable binding characteristics may still be absent from the cell-surface immunopeptidome because its source protein is not expressed or degraded, the peptide is not generated or transported efficiently, or it is not detected experimentally. Binding scores therefore do not directly encode presentation-evidence preferences among tissues that share the same MHC restriction.

\subsection{Antigen Presentation Prediction}

Eluted-ligand data enable models to learn a broader presentation signal than binding measurements alone \cite{reynisson2020netmhcpan,odonnell2020mhcflurry}. General presentation predictors combine peptide sequence, MHC context, and, in some cases, antigen-processing or source-protein features to distinguish observed ligands from background peptides. Systems such as BigMHC \cite{albert2023bigmhc} illustrate how large immunopeptidomic collections and modern sequence models can improve ligand presentation prediction and candidate ranking. Their usual target, however, is whether a peptide is likely to be presented for an allele, sample, or pooled presentation setting. This target differs from the relative question studied here: given a fixed tissue--MHC task, which of two peptides from the same parent protein has stronger tissue-conditioned presentation evidence? Consequently, a high general-presentation score need not imply a tissue-specific preference, and direct comparisons require a benchmark whose supervision reflects that distinction.

\subsection{Tissue- and Sample-Context Immunopeptidomics}

Multi-tissue immunopeptidomic atlases have revealed that ligand repertoires contain both broadly shared peptides and tissue-associated components \cite{marcu2021atlas,schuster2018mouse}. MHCflurry 2.0 incorporates antigen-processing features but does not use tissue identity \cite{odonnell2020mhcflurry}. More recent systems support multi-allelic deconvolution and ligand analysis, as in ImmuneApp \cite{xu2024immuneapp}, or combine peptide, MHC, flanking sequence, and optional gene-expression features, as in HLApollo \cite{thrift2024hlapollo}. These approaches demonstrate that presentation is context dependent, but they generally do not define a tissue--MHC combination as the supervised task or match every positive and pseudo-negative on both MHC restriction and parent protein. Our formulation is therefore complementary to expression- and sample-aware prediction: it isolates a task-conditioned ranking problem that can later be combined with molecular context.

\subsection{Multi-task Learning for Biological Sequences}

Biological prediction tasks often exhibit unequal sample sizes, related labels, and heterogeneous sequence patterns, making multi-task learning a natural strategy for statistical sharing \cite{caruana1997multitask,ruder2017overview}. A shared encoder with task-specific heads can transfer common sequence features while retaining task-level decision boundaries. More structured architectures, including multi-gate mixture-of-experts \cite{ma2018mmoe} and progressive layered extraction \cite{tang2020ple}, allow tasks or task groups to select among shared and specialized representations. Auxiliary objectives, task grouping, and gradient-balancing methods provide additional mechanisms for improving transfer or limiting negative interference. Nevertheless, unrestricted sharing may obscure allele-specific motifs, whereas fully independent models waste information and are poorly suited to long-tailed tasks. The TissuePMHC human model addresses this trade-off by combining global supervision with HLA-specific sharing and task-specific prediction, using a position-preserving multi-scale peptide representation designed for the fixed-length MHC-I setting.

\subsection{Relation to This Work}

Unlike methods that predict general peptide--MHC binding or presentation, or that use tissue expression only as an additional input, our work directly defines each tissue--MHC combination as a target task. Supervision is constructed from peptide pairs matched by MHC restriction and parent UniProt protein, so the benchmark focuses on within-protein, task-conditioned presentation-evidence preference rather than absolute presentation probability. We further distinguish the task definition from the model used to solve it: the human benchmark motivates the dual-branch TissuePMHC human model, whereas the smaller mouse benchmark uses a separately trained structured multi-task model to test whether task learnability and the benefit of structured sharing are reproducible at the benchmark level. Evaluation under both standard pair-disjoint and connected-component peptide-disjoint protocols explicitly separates primary closed-task performance from seen-task, unseen-peptide behavior. Frozen MHCflurry and NetMHCpan scores are evaluated on the same pairs as general-signal controls rather than as models of tissue context. Relative to the closest context-aware presentation models \cite{odonnell2020mhcflurry,albert2023bigmhc}, the TissuePMHC benchmark differs along three axes: the prediction target is tissue--MHC preference rather than absolute presentation, every pseudo-negative is matched to a positive on both MHC restriction and parent protein, and evaluation explicitly separates pair-disjoint performance from peptide-disjoint generalization.

Table~\ref{tab:closest-work} makes these distinctions explicit. Negative construction varies across implementations and datasets; the entries therefore summarize the published task design rather than asserting identical data provenance.

\begin{table}[tbp]
\centering
\caption{Task-level comparison with representative peptide--MHC methods. ``Context'' denotes information beyond peptide and MHC used by the published method.}
\label{tab:closest-work}
\scriptsize
\setlength{\tabcolsep}{2pt}
\begin{tabular}{@{}>{\raggedright\arraybackslash}p{1.8cm}>{\raggedright\arraybackslash}p{2.5cm}>{\raggedright\arraybackslash}p{2.3cm}>{\raggedright\arraybackslash}p{2.6cm}@{}}
\toprule
Method & Prediction target & Context & Negative/evaluation emphasis \\
\midrule
NetMHCpan 4.1 \cite{reynisson2020netmhcpan} & Binding and eluted-ligand presentation & MHC pseudo-sequence; no tissue task & Pan-allelic general prediction \\
MHCflurry 2.0 \cite{odonnell2020mhcflurry} & Binding and presentation & Antigen-processing features; no tissue identity & Pan-allelic general prediction \\
BigMHC \cite{albert2023bigmhc} & Presentation and immunogenicity & Peptide and MHC & Large-scale general prediction \\
ImmuneApp \cite{xu2024immuneapp} & HLA-I epitope prediction and deconvolution & Mono- and multi-allelic ligand data & General/sample-level evaluation \\
HLApollo \cite{thrift2024hlapollo} & Peptide presentation & MHC, peptide flanks, optional gene expression & Negative-set switching; pan-allelic evaluation \\
TissuePMHC benchmark & Tissue--MHC presentation preference & Tissue--MHC task identity & MHC- and parent-matched pairs; pair- and peptide-disjoint evaluation \\
\bottomrule
\end{tabular}
\end{table}
